\documentclass[12pt,letterpaper]{article}
\usepackage{graphicx,textcomp}
\usepackage{natbib}
\usepackage{setspace}
\usepackage{fullpage}
\usepackage{color}
\usepackage[reqno]{amsmath}
\usepackage{amsthm}
\usepackage{fancyvrb}
\usepackage{amssymb,enumerate}
\usepackage[all]{xy}
\usepackage{endnotes}
\usepackage{lscape}
\newtheorem{com}{Comment}
\usepackage{float}
\usepackage{hyperref}
\newtheorem{lem} {Lemma}
\newtheorem{prop}{Proposition}
\newtheorem{thm}{Theorem}
\newtheorem{defn}{Definition}
\newtheorem{cor}{Corollary}
\newtheorem{obs}{Observation}
\usepackage[compact]{titlesec}
\usepackage{dcolumn}
\usepackage{tikz}
\usetikzlibrary{arrows}
\usepackage{multirow}
\usepackage{xcolor}
\newcolumntype{.}{D{.}{.}{-1}}
\newcolumntype{d}[1]{D{.}{.}{#1}}
\definecolor{light-gray}{gray}{0.65}
\usepackage{url}
\usepackage{listings}
\usepackage{color}

\definecolor{codegreen}{rgb}{0,0.6,0}
\definecolor{codegray}{rgb}{0.5,0.5,0.5}
\definecolor{codepurple}{rgb}{0.58,0,0.82}
\definecolor{backcolour}{rgb}{0.95,0.95,0.92}

\lstdefinestyle{mystyle}{
	backgroundcolor=\color{backcolour},   
	commentstyle=\color{codegreen},
	keywordstyle=\color{magenta},
	numberstyle=\tiny\color{codegray},
	stringstyle=\color{codepurple},
	basicstyle=\footnotesize,
	breakatwhitespace=false,         
	breaklines=true,                 
	captionpos=b,                    
	keepspaces=true,                 
	numbers=left,                    
	numbersep=5pt,                  
	showspaces=false,                
	showstringspaces=false,
	showtabs=false,                  
	tabsize=2
}
\lstset{style=mystyle}
\newcommand{\Sref}[1]{Section~\ref{#1}}
\newtheorem{hyp}{Hypothesis}

\title{Problem Set 3}
\date{Due: March 25, 2026}
\author{Applied Stats II}


\begin{document}
	\maketitle
	\section*{Instructions}
	\begin{itemize}
		\item Please show your work! You may lose points by simply writing in the answer. If the problem requires you to execute commands in \texttt{R}, please include the code you used to get your answers. Please also include the \texttt{.R} file that contains your code. If you are not sure if work needs to be shown for a particular problem, please ask.
		\item Your homework should be submitted electronically on GitHub in \texttt{.pdf} form.
		\item This problem set is due before 23:59 on Wednesday March 25, 2026. No late assignments will be accepted.
	\end{itemize}
	
	\vspace{.25cm}
	\section*{Question 1}
	\vspace{.25cm}
	\noindent We are interested in how governments' management of public resources impacts economic prosperity. Our data come from \href{https://www.researchgate.net/profile/Adam_Przeworski/publication/240357392_Classifying_Political_Regimes/links/0deec532194849aefa000000/Classifying-Political-Regimes.pdf}{Alvarez, Cheibub, Limongi, and Przeworski (1996)} and is labelled \texttt{gdpChange.csv} on GitHub. The dataset covers 135 countries observed between 1950 or the year of independence or the first year forwhich data on economic growth are available ("entry year"), and 1990 or the last year for which data on economic growth are available ("exit year"). The unit of analysis is a particular country during a particular year, for a total $>$ 3,500 observations. 
	
	\begin{itemize}
		\item
		Response variable: 
		\begin{itemize}
			\item \texttt{GDPWdiff}: Difference in GDP between year $t$ and $t-1$. Possible categories include: "positive", "negative", or "no change"
		\end{itemize}
		\item
		Explanatory variables: 
		\begin{itemize}
			\item
			\texttt{REG}: 1=Democracy; 0=Non-Democracy
			\item
			\texttt{OIL}: 1=if the average ratio of fuel exports to total exports in 1984-86 exceeded 50\%; 0= otherwise
		\end{itemize}
		
	\end{itemize}
	\newpage
	\noindent Please answer the following questions:
	
	\begin{enumerate}
		\item Construct and interpret an unordered multinomial logit with \texttt{GDPWdiff} as the output and ``no change" as the reference category, including the estimated cutoff points and coefficients.
		
		\lstinputlisting[
		language=R,
		firstline=47,
		lastline=59
		]{PS3_SK.R}
		
		\begin{verbatim}
			# weights:  12 (6 variable)
			initial  value 4087.936326 
			iter  10 value 2340.076844
			final  value 2339.385155 
			convergedtent...
			
			Coefficients:
												(Intercept)      REG      OIL
			negative    3.805370 1.379282 4.783968
			positive    4.533759 1.769007 4.576321
			
			Std. Errors:
												(Intercept)       REG      OIL
			negative   0.2706832 0.7686958 6.885366
			positive   0.2692006 0.7670366 6.885097
			
			> p_values_multi
												(Intercept)        REG       OIL
			negative           0 0.07276308 0.4871792
			positive           0 0.02109459 0.5062612
		\end{verbatim}
		
		\vspace{0.2cm}
		We estimated an unordered multinomial logit model with ``no change" as the reference category.
		\begin{itemize}
			\item \textbf{Comparing ``negative" to ``no change":} The estimated intercept for the ‘negative’ versus ‘no change’ comparison is $3.805$. Being a democracy (\texttt{REG}=1) increases the log-odds of experiencing negative GDP growth (compared to no change) by $1.379$, holding \texttt{OIL} constant ($p = 0.073$).
			\item \textbf{Comparing ``positive" to ``no change":} The estimated intercept for the ‘positive’ versus ‘no change’ comparison is $4.534$. Being a democracy (\texttt{REG}=1) significantly increases the log-odds of experiencing positive GDP growth (compared to no change) by $1.769$, holding \texttt{OIL} constant ($p = 0.021$).
			
			\item The OIL coefficient is positive in both comparisons, suggesting oil-exporting countries may have higher log-odds of being in either the negative or positive GDP-change category relative to ‘no change,’ but these estimates are not statistically significant.
		\end{itemize}
		Overall, democracy is clearly associated with higher log-odds of positive GDP change relative to no change. For negative GDP change relative to no change, the coefficient is also positive, but the evidence is weaker and not statistically significant at the 5\% level.

		\vspace{0.2cm}
		
		\item Construct and interpret an ordered multinomial logit with \texttt{GDPWdiff} as the outcome variable, including the estimated cutoff points and coefficients.
		
		\lstinputlisting[
		language=R,
		firstline=61,
		lastline=72
		]{PS3_SK.R}
		
		\begin{verbatim}
																											Value Std. Error    t value      p value
			REG                 0.3984834 0.07518479   5.300054 1.157687e-07
			OIL                -0.1987177 0.11571713  -1.717271 8.592967e-02
			negative|no change -0.7311784 0.04760375 -15.359680 0.000000e+00
			no change|positive -0.7104851 0.04750680 -14.955440 0.000000e+00
		\end{verbatim}
		
		We estimated an ordered multinomial logit model with categories ordered as $\text{negative} < \text{no change} < \text{positive}$.
		\begin{itemize}
			\item \textbf{Cutoff points (Intercepts):} The estimated thresholds are -0.731 for the negative/no-change split and -0.710 for the no-change/positive split. The coefficient on REG is 0.398, meaning democracy increases the log-odds of being in a higher GDP-change category, holding OIL constant, and this effect is highly statistically significant (p < 0.001).
			\item \textbf{Coefficients:} The coefficient on OIL is -0.199, suggesting a lower likelihood of being in a higher GDP-change category, but this estimate is not statistically significant at the 5\% level (p = 0.086).
		\end{itemize}
		
	\end{enumerate}
	
	\section*{Question 2} 
	\vspace{.25cm}
	
	\noindent Consider the data set \texttt{MexicoMuniData.csv}, which includes municipal-level information from Mexico. The outcome of interest is the number of times the winning PAN presidential candidate in 2006 (\texttt{PAN.visits.06}) visited a district leading up to the 2009 federal elections, which is a count. Our main predictor of interest is whether the district was highly contested, or whether it was not (the PAN or their opponents have electoral security) in the previous federal elections during 2000 (\texttt{competitive.district}), which is binary (1=close/swing district, 0="safe seat"). We also include \texttt{marginality.06} (a measure of poverty) and \texttt{PAN.governor.06} (a dummy for whether the state has a PAN-affiliated governor) as additional control variables. 
	
	\begin{enumerate}
		\item [(a)]
		Run a Poisson regression because the outcome is a count variable. Is there evidence that PAN presidential candidates visit swing districts more? Provide a test statistic and p-value.
		
		\lstinputlisting[
		language=R,
		firstline=82,
		lastline=92
		]{PS3_SK.R}
		
		\begin{verbatim}
			Coefficients:
																								Estimate Std. Error z value Pr(>|z|)    
			(Intercept)          -3.81023    0.22209 -17.156   <2e-16 ***
			competitive.district -0.08135    0.17069  -0.477   0.6336    
			marginality.06       -2.08014    0.11734 -17.728   <2e-16 ***
			PAN.governor.06      -0.31158    0.16673  -1.869   0.0617 .  
			---
			Signif. codes:  0 ‘***’ 0.001 ‘**’ 0.01 ‘*’ 0.05 ‘.’ 0.1 ‘ ’ 1
			
			(Dispersion parameter for poisson family taken to be 1)
			
			Null deviance: 1473.87  on 2406  degrees of freedom
			Residual deviance:  991.25  on 2403  degrees of freedom
			AIC: 1299.2
			
			Number of Fisher Scoring iterations: 7
		\end{verbatim}
		
		\begin{verbatim}
			> summary(poisson_model)$coefficients["competitive.district", ]
						Estimate  Std. Error     z value    Pr(>|z|) 
			-0.08135181  0.17068819 -0.47661062  0.63363942 
			
			> anova(poisson_reduced, poisson_model, test="LRT")
			Analysis of Deviance Table
					Resid. Df Resid. Dev Df Deviance Pr(>Chi)
			1      2404     991.48                     
			2      2403     991.25  1  0.22377   0.6362
		\end{verbatim}
		
		\vspace{0.2cm}
		\textbf{No. There is no statistically significant evidence that PAN presidential candidates visited swing districts more often.} In fact, the estimated coefficient on competitive.district is negative (-0.081), implying a slightly lower expected visit count in competitive districts, but this effect is substantively small and statistically insignificant (Wald z = -0.477, p = 0.634; LRT p = 0.636).
		
		\item [(b)]
		Interpret the \texttt{marginality.06} and \texttt{PAN.governor.06} coefficients.
		
		\lstinputlisting[
		language=R,
		firstline=94,
		lastline=96
		]{PS3_SK.R}
		
		\begin{verbatim}
			
(Intercept) competitive.district       marginality.06      PAN.governor.06 
						0.022                0.922                0.125                0.732 
		\end{verbatim}
		
		\vspace{0.2cm}
		\textit{(Note: Exponentiating the Poisson coefficients yields Incidence Rate Ratios).}
		\begin{itemize}
			\item \textbf{\texttt{marginality.06}:} For a one-unit increase in the poverty level, the expected number of PAN presidential visits changes by a \textbf{multiplicative factor of 0.125} ($e^{-2.080}$), holding other variables constant. In other words, visits drastically decrease as poverty increases.
			\item \textbf{\texttt{PAN.governor.06}:} The incidence rate ratio for PAN.governor.06 is 0.732, suggesting that districts in states with a PAN-affiliated governor are expected to receive about 26.8\% fewer visits than districts in states without a PAN governor, holding other variables constant. However, this estimate is not statistically significant at the 5\% level (p = 0.0617).
		\end{itemize}
		
		
		\item [(c)]
		Provide the estimated mean number of visits from the winning PAN presidential candidate for a hypothetical district that was competitive (\texttt{competitive.district}=1), had an average poverty level (\texttt{marginality.06} = 0), and a PAN governor (\texttt{PAN.governor.06}=1).
		
		\lstinputlisting[
		language=R,
		firstline=98,
		lastline=100
		]{PS3_SK.R}
		
		\begin{verbatim}
			1 
			0.01
		\end{verbatim}
		
		\lstinputlisting[
		language=R,
		firstline=102,
		lastline=104
		]{PS3_SK.R}
		
		\begin{verbatim}
(Intercept) 
			0.01 
		\end{verbatim}
		
		\vspace{0.2cm}
		For this hypothetical district, the estimated mean number of visits from the winning PAN candidate is \textbf{0.01 visits} on average. \\
		\textit{(Calculated using the inverse log link: $\exp(-3.810 - 0.081(1) - 2.080(0) - 0.312(1)) \approx 0.01$)}.
		
	\end{enumerate}
	
\end{document}
